\section{Dataset}
\label{sec:dataset}


\subsection{Descriptor Families}

\textbf{Partial charges.}
Six independent atomic partial-charge schemes are computed per structure.
Hirshfeld \citep{hirshfeld1977}, CM5 \citep{cm5}, ADCH \citep{adch}, and Becke
\citep{becke1988} charges are produced by Multiwfn \citep{multiwfn} from the DFT electron density;
Mulliken and L\"owdin charges are extracted directly from the ORCA output \citep{orca5}.
Hirshfeld charges partition the electron density by the stockholder principle: each
atom claims a share of the density at every point proportional to its free-atom
promolecular density, yielding compact, chemically intuitive atomic populations.
CM5 applies a pairwise connectivity-based correction to Hirshfeld charges to
reproduce experimental dipole moments across diverse environments \citep{cm5}.
Becke charges integrate the molecular density over smooth, fuzzy atomic basins
defined by Becke's numerical weight functions \citep{becke1988}, providing a
partition that is independent of the basis set and free of the basis-function
artifacts that affect Mulliken and L\"owdin populations.
ADCH further corrects Hirshfeld charges for atomic dipole contributions to reduce
basis-set dependence \citep{adch}.

\textbf{Bond orders.}
Four schemes are universal across the corpus: Mayer \citep{mayer1983} and
L\"owdin bond orders from ORCA, fuzzy bond orders from Multiwfn
\citep{multiwfn,becke1988}, and QTAIM bond presence inferred from bond critical points.
Mayer bond orders derive from the density matrix and capture bond multiplicity
beyond formal valence, including partial and dative character.
L\"owdin bond orders use a symmetrically orthogonalized basis, reducing
the basis-set sensitivity that affects Mulliken-derived indices.
Fuzzy bond orders integrate the exchange-correlation hole over Becke fuzzy-atom
basins, providing a real-space, basis-set-independent measure of covalent
bond character. QTAIM bond presence is a binary signal at each bond critical
point; the corresponding density and Laplacian scalars are stored under QTAIM topology.

\textbf{QTAIM topology.}
Critical points of the electron density are located via Multiwfn following Bader's
Quantum Theory of Atoms in Molecules \citep{bader1990}.
Each bond critical point carries 26 scalar fields ($\rho$, $\nabla^2\rho$,
ellipticity, kinetic/potential energy densities); ring and cage CPs, non-nuclear
attractors, and atomic basin integrals are stored where present.

\textbf{Fuzzy and surface descriptors.}
Fuzzy atom analysis partitions molecular properties via smooth, overlapping atomic
weight functions, yielding basis-set-independent atomic populations, valences, and
bond orders~\citep{mayersalvador2004}.
Per-atom Becke and Hirshfeld fuzzy electron-density integrations \citep{becke1988}
are computed via Multiwfn for every structure (with spin-density counterparts for
open-shell records).
These provide atomic electron and spin populations that are independent of the
one-electron basis. Molecular-level ALIE (Average Local Ionization Energy) 
surface properties, including surface volume, area, extrema, polarity index, 
and internal charge separation are additionally extracted per molecule, 
encoding site-specific electrophilicity and reactivity complementary 
to the atomic-level descriptors.

\textbf{ORCA globals.}
Total SCF energy, HOMO/LUMO orbital energies, molecular dipole, nuclear gradient, and
SCF convergence metadata are parsed from the ORCA output file \citep{orca5,ORCA}.


\subsection{Storage Layout}
\label{sec:storage}

Each split folder contains eight LMDB files keyed on a single string identifier
$k = \texttt{relpath(folder, root).replace("/", "\_\_")}$. Keys for each structure
are shared across all lmdbs and they each contain different datatypes:
\begin{itemize}
\item \texttt{structure.lmdb}: pymatgen \texttt{Molecule} (positions, elements), total spin and charge, RDKit-derived bond list, and external identifiers.
\item \texttt{charge.lmdb}: per-atom charges nested by scheme (\texttt{hirshfeld}, \texttt{cm5}, \texttt{adch}, \texttt{becke}, \texttt{mulliken\_orca}, \texttt{loewdin\_orca}); \texttt{hirshfeld}/\texttt{cm5} also carry a total dipole, \texttt{adch}/\texttt{becke} carry total and per-atom dipoles.
\item \texttt{qtaim.lmdb}: per-structure dict of critical points keyed by atom ID for nuclear CPs or by atom-pair for bond CPs; 26 scalar fields per CP including $\rho$, $\nabla^{2}\rho$, ellipticity, kinetic and potential energy densities.
\item \texttt{bond.lmdb}: bond-order schemes (\texttt{fuzzy\_bond}, \texttt{mayer\_orca}, \texttt{loewdin\_orca}) keyed by directed atom-pair strings \texttt{"\{i\}\_\{El\_i\}\_to\_\{j\}\_\{El\_j\}"}; the fourth scheme, QTAIM bond presence, is recoverable from \texttt{qtaim.lmdb} BCP keys.
\item \texttt{fuzzy.lmdb}: per-atom Multiwfn fuzzy integrations (\texttt{becke\_fuzzy\_density}, \texttt{hirsh\_fuzzy\_density}; spin counterparts for open-shell records).
\item \texttt{other.lmdb}: molecular descriptors from ESP and ALIE surface analysis (volume, surface area, min/max, signed mean and variance, polarity index, internal charge separation) plus four polarity scalars (\texttt{mpp\_full}, \texttt{sdp\_full}, \texttt{mpp\_heavy}, \texttt{sdp\_heavy}).
\item \texttt{orca.lmdb}: ORCA-derived globals (total energy, orbital energies, dipole, gradient, SCF metadata), filtered via \texttt{DEFAULT\_ORCA\_FILTER} when consumed by the converter.
\item \texttt{timings.lmdb}: per-job step durations; provenance only.
\end{itemize}


\subsection{Data Quality and Per-Vertical Breakdown}
\label{sec:dataset-quality}

Calculations and conversions to LMDBs are highly stable with five of the
seven analytic LMDBs showing no corrupt or incomplete data storage. 
For \texttt{qtaim} and \texttt{orca} we have 25 (0.0006\%) and 1{,}607 (0.0403\%)
failed calculations; the ORCA failures concentrate on physically
hard cases (solvated biomolecules, scaled-separation electrolytes,
open-shell TM redox), and geometry plus Multiwfn-side descriptors stay
valid for those records. The failed QTAIM records were instances where
no BCPs were detected or where the Poincare-Hopf formula was not obeyed.
Cross-LMDB key alignment holds to within four
records on a 4M-key index, so any descriptor combination joins on
the full corpus. Each of the six charge schemes (Hirshfeld, ADCH, CM5,
Becke, Mulliken, L\"owdin), four bond-order schemes (fuzzy,
Mayer, L\"owdin, QTAIM), full QTAIM topology, and the standard ORCA globals
reaches $\geq$99\% on every one of the 34 verticals; in other words, the same
profile is available for a noble-gas dimer, a 200-atom solvated protein
pocket, and a transition-metal redox complex. Full audit
table, coverage heatmap, and per-vertical record counts:
App.~\ref{app:lmdb_audit} and Table~\ref{tab:t1_census}.
